% Pursuing Stacks (cotes 134-2 to 134-8): archive cote 134-6, site page 25 (PDF page 26), typescript page 344.
% typescript. Draft compared: ps5.tex:2596-2667.
% AI-transcribed by Claude Opus 5.5 (claude-opus-5-5, adaptive thinking, effort high), Message Batches API, batch ps3 (msgbatch_01XdFmuVVBAcgc7EFxdtCqsj), 2026-11-12; not yet checked by eye.
\scanpage{134-6/25}
\mspage{344}

Still, we will call a category A satisfying the condition (Bas 4)
with respect to the embedding
\[ (1) \qquad A \hand{\hookrightarrow} A^{\wedge} \]
a \uline{precontractor} , as we'll expect something more still from a contrac-%
tor, which \xout{\unsure{ma}} will be automatically satisfied in the particular case
when A is stable under finite products. Roughly speaking we want to
have a satisfactory relation between contractibility and asphericity
in $A^{\wedge}$ - we'll make this more precise below. For the time being, \xout{xxxx}
let's dwell just a little more on the notion of a precontractor.

A second thought about contractors, coming alongside with the
first, is that for any full embedding of A into a larger category
\[ (2) \qquad f\colon A \hand{\longrightarrow} M \ , \qquad M \text{ stable under finite products ,} \]
f(A) should generate in M a contractibility structure. In the particu-%
lar case when A is stable under finite products (and hence the notion
of a precontractor\hand{,} already defined, coincides with the notion of a
contractor), this is indeed so provided A is \deleted{indeed} a (pre)contractor,
and moreover f commutes to finite products. When A is just assumed to
be a \xout{contractor} precontractor (without an assumption about stability of
A under products), we'll assume in compensation that M is stable under
small direct limits, which allows to take the canonical extension $f_{!}$
of f to $A^{\wedge}$ in a way commut\xout{\unsure{a}}ing to \xout{products} direct limits
\[ (3) \qquad f_{!} : A^{\wedge} \hand{\longrightarrow} M \ , \]
and we can now state: \uline{if} $f_{!}$ \uline{commutes to finite products} (cf. prop.1 ,
a) p. 281), \uline{then} f(A) \uline{generates a contractibility structure in} M.
This statement is true even without assuming that the functor f is ful-%
ly faithful (and follows immediately from the criterion (Bas 4) of
p. 118); however, in the particular case when f is fully faithful,
we have a handy criterion \inserted[typed]{(prop.2, p.283)} for $f_{!}$ to commute to finite products, namely
that f(A) be a strictly generating subcategory of M, or equivalently,
that the functor
\[ \hand{(4)} \qquad f^{*}\colon M \hand{\to} A^{\wedge} \]
(right adjoint to $f_{!}$) be fully faithful. In this case, we may identi-%
fy M (up to equivalence) to a full subcategory of $A^{\wedge}$ containing A ,
and the fact that f(A) generates a contractibility structure in M
follows immediately directly (without having to rely on existence of
direct limits in M, nor even existence of $f_{!}$) . To sum up:

\margin{25}
